\Chapter{12}

\Verse{1}
{
	\hJ{}
}
{
	\eJ{
		And YHWH said to Abram, “Go from your land and from your birthplace and
from your father’s house to the land that I’ll show you.\\
	}
}
{
}

\Verse{2}
{
	\hJ{}
}
{
	\eJ{
		And I’ll make you into a big nation, and I’ll bless you and make your
name great. And be a blessing!\\
	}
}
{
}

\Verse{3}
{
	\hJ{}
}
{
	\eJ{
		And I’ll bless those who bless you, and those who affront you I’ll
curse. And all the families of the earth will be blessed through you.”\\
	}
}
{
}

\Verse{4}
{
	\hJ{}
}
{
	\eJ{
		And Abram went as YHWH had spoken to him, and Lot went with him. And
Abram was seventy-five years old when he went out from Haran.\\
	}
}
{
}

\Verse{5}
{
	\hP{}
}
{
	\eP{
		And Abram took Sarai, his wife, and Lot, his brother’s son, and all
their property that they had accumulated and the persons whom they had
gotten in Haran; and they went out to go to the land of Canaan.

And they came to the land of Canaan.\\
	}
}
{
}

\Verse{6}
{
	\hJ{}
}
{
	\eJ{
		And Abram passed through the land as far as the place of Shechem, as
far as the oak of Moreh. And the Canaanite was in the land then.\\
	}
}
{
}

\Verse{7}
{
	\hJ{}
}
{
	\eJ{
		And YHWH appeared to Abram and said, “I’ll give this land to your
seed.” And he built an altar there to YHWH who had appeared to him.\\
	}
}
{
}

\Verse{8}
{
	\hJ{}
}
{
	\eJ{
		And he moved on from there to the hill country east of Beth-El and
pitched his tent—Beth-El to the west and Ai to the east—and he built an
altar there to YHWH and invoked the name YHWH.\\
	}
}
{
}

\Verse{9}
{
	\hJ{}
}
{
	\eJ{
		And Abram traveled on, going to the Negeb.\\
	}
}
{
}

\Verse{10}
{
	\hJ{}
}
{
	\eJ{
		And there was a famine in the land, and Abram went down to Egypt to
reside there because the famine was heavy in the land.\\
	}
}
{
}

\Verse{11}
{
	\hJ{}
}
{
	\eJ{
		And it was when he was close to coming to Egypt: and he said to Sarai,
his wife, “Here, I know that you’re a beautiful woman.

12:11–13. Here, I know. Abraham’s words to Sarah contain the first two
occurrences of the Hebrew word na’. It is sometimes rendered “now” in
English and sometimes rendered “please.” But neither is quite right. It
is an untranslatable particle that is a sign of polite speech. Its most
striking occurrence is in God’s words to Abraham in the story of the
binding of Isaac: “Take na’ your son … and offer him” (Gen 22:2). I do
not translate it with any English word, because it is closer to the
Hebrew to use no word than to use such words as “please” or “now.”\\
	}
}
{
}

\Verse{12}
{
	\hJ{}
}
{
	\eJ{
		And it will be when the Egyptians see you that they’ll say, ‘This is
his wife,’ and they’ll kill me and keep you alive.\\
	}
}
{
}

\Verse{13}
{
	\hJ{}
}
{
	\eJ{
		Say you’re my sister so it will be good for me on your account and
I’ll stay alive because of you.”\\
	}
}
{
}

\Verse{14}
{
	\hJ{}
}
{
	\eJ{
		And it was as Abram came to Egypt, and the Egyptians saw the woman,
that she was very beautiful.\\
	}
}
{
}

\Verse{15}
{
	\hJ{}
}
{
	\eJ{
		And Pharaoh’s officers saw her and praised her to Pharaoh, and the
woman was taken to Pharaoh’s house.\\
	}
}
{
}

\Verse{16}
{
	\hJ{}
}
{
	\eJ{
		And he was good to Abram on her account, and he had a flock and oxen
and he-asses and servants and maids and she-asses and camels.\\
	}
}
{
}

\Verse{17}
{
	\hJ{}
}
{
	\eJ{
		And YHWH plagued Pharaoh and his house, big plagues, over the matter
of Sarai, Abram’s wife.\\
	}
}
{
}

\Verse{18}
{
	\hJ{}
}
{
	\eJ{
		And Pharaoh called Abram and said, “What is this you’ve done to me?
Why didn’t you tell me that she was your wife?\\
	}
}
{
}

\Verse{19}
{
	\hJ{}
}
{
	\eJ{
		Why did you say, ‘She’s my sister,’ so I took her for a wife for
myself! And now, here’s your wife. Take her and go.”\\
	}
}
{
}

\Verse{20}
{
	\hJ{}
}
{
	\eJ{
		And Pharaoh commanded people over him, and they let him and his wife
and all he had go.\\
	}
}
{
}

